\begin{center}
{\Large \bfseries \textcolor{lyred}{第三章 微分中值定理与导数的应用}}
\end{center}

\vspace{6pt}
{\large \bfseries \textcolor{lyred}{第3.1 节 微分中值定理}}

\vspace{4pt}
{\bfseries \textcolor{lyred}{一.罗尔定理}}

\textcolor{lyred}{费马引理}.设在
(
)
U x
内
()
(
)
f x
f x
\le 
或
()
(
)
f x
f x
\ge 
,若
()
f x 在
0x 处可导,则 
(
)
f
x
'
=
,即
0x 为
()
f x 的\textcolor{lyred}{驻点}. 
\textcolor{lyred}{定理}.设
()
f x 在[
]
,a b 上连续,在(
)
,a b 内可导,若
()
()
f a
f b
=
,则存在
(
)
,a b
\psi \in 
, 
使得
()
f
\psi 
'
=
,其中\psi 称为\textcolor{lyred}{中值}. 
\textcolor{lyred}{注}.我们约定本节的函数均满足:在[
]
,a b 上连续,(
)
,a b 内可导. 
\textcolor{lyred}{例}.设
()
()
f
f
+
=
,
()
f
=
,证明:
(
)
0,2
\psi 
\exists \in 
,使得
()
f
\psi 
'
=
. 
证.
[
]
0,1
\eta 
\exists \in 
,使得
()
()
()
f
f
f \eta 
+
=
=
(加权平均值),证毕. 
\textcolor{lyred}{例}.设
()
()
f
f
>
,
()
()
f
f
>
,证明:存在
(
)
0,2
\psi \in 
,使得
()
f
\psi 
'
=
. 
证.不妨设
()
()
()
f
f
f
>
>
,则
(
)
1,2
\eta 
\exists \in 
,使得
()
()
f
f
\eta =
,即得; 
或者,
(
)
0,2
\psi 
\exists \in 
,使得
()
()
min
x
f
f x
\psi 
\le \le 
=
,由费马引理,即得,证毕. 
\textcolor{lyred}{例}.设
()
()
f
f
=
=
,
()1
f
=
,证明:存在
(
)
0,2
\psi \in 
,使得
()
f
\psi 
'
=. 
证.令
()
()
F x
f x
x
=
-
,则
()
F
=
,
()1
F
=,
()
F
=-,故
(
)
0,2
\psi 
\exists \in 
,使得 
()
F \psi 为最大值,于是
()
F \psi 
'
=
,证毕. 
\textcolor{lyred}{例}.设曲线
()
y
f x
=
与某一条不平行于y 轴的直线相交于三个点,证明:存在\psi , 
使得
()
f
\psi 
''
=
. 
证.设直线方程为y
ax
b
=
+
,令
()
()
F x
f x
ax
b
=
-
-,则
()
F x 有三个不同零点,故 
存在\psi ,使得
()
F
\psi 
''
=
,证毕. 
\textcolor{lyred}{注}.若曲线
()
y
f x
=
与一条抛物线交于四个点,则存在\psi ,使得
()
f
\psi 
'''
=
. 
\textcolor{lyred}{注}.(1)若
()
f x 有
n +个零点,则
()
f
x
'
至少有n 个零点,
()
f
x
''
至少有
n -个, 
…,
()()
n
f
x 至少有1个零点. 
(2)若
()()
n
f
x \neq 
,例如n 次多项式,则
()
f x 至多有n 个零点. 
\textcolor{lyred}{例}.设
()
(
)(
)
(
)
1 2
f x
x x
x
nx
=
-
-
-
\Lambda 
,求
()
f
x
''
的实零点个数. 
解.
()
f x 有
n +个零点,
()
f
x
''
至少有
n -个零点,由于
()
f
x
''
为
n -次多项式, 
至多有
n -个零点,故恰有
n -个零点. 
\textcolor{lyred}{例}.证明:方程
x
x
-
=有且仅有3个实根. 
证.设
()
x
f x
x
=
-
-,则
()
()
f
f
=
=
,又
()
f
=-,
()
f
=
,故 
有3个实根,同时
()(
)
ln 2
x
x
f
x
'''
'''
=
=
\cdots 
\neq 
,证毕. 
\textcolor{lyred}{例}.设
()1
f
=
,
()
()
F x
x f x
=
,证明:存在
(
)
0,1
\psi \in 
,使得
()
F
\psi 
''
=
. 
证.
()
()
(
)
0,1
F
F
\eta 
=
=
\Rightarrow \exists \in 
,使得
()
F \eta 
'
=
,又
()
F'
=
,证毕. 
\textcolor{lyred}{例}.设
()
()
f
f
=
=
,
()
()
F x
xf x
=
,证明:存在
(
)
0,1
\psi \in 
,使得
()
F
\psi 
''
=
. 
证.
()
()
(
)
0,1
F
F
\eta 
=
=
\Rightarrow \exists \in 
,使得
()
F \eta 
'
=
,又
()
F'
=
,证毕. 
\textcolor{lyred}{例}.设
()1
f
=
,证明:存在
(
)
0,1
\psi \in 
,使得
()
()
f
f
\psi 
\psi 
\psi 
'
+
=
. 
证.令
()
()
F x
xf x
=
,则
()
()
F
F
=
=
,而
()
()
()
F
x
f x
xf
x
'
'
=
+
,证毕. 
\textcolor{lyred}{例}.设
()
()
f a
f b
=
=
,证明:存在
(
)
,a b
\psi \in 
,使得
()
()
f
f
\psi 
\psi 
'
+
=
. 
证.令
()
()
x
F x
e f x
=
,则
()
()
F a
F b
=
=
,而
()
()
()
x
F
x
e
f x
f
x
'
'
=
+



,证毕. 
\vspace{4pt}
{\bfseries \textcolor{lyred}{二.拉格朗日中值定理}}

\textcolor{lyred}{定理}.设
()
f x 在[
]
,a b 上连续,(
)
,a b 内可导,则存在
(
)
,a b
\psi \in 
,使得 
()
()
()
f b
f a
f
b
a
\psi 
-
'
=
-
. 
\textcolor{lyred}{推论(增量公式)}.
(
)
()
(
)
y
f x
x
f x
f
x
x
x
\Delta 
\Delta 
\theta \Delta 
\Delta 
'
=
+
-
=
+
,其中0
\theta 
<
<. 
\textcolor{lyred}{推论}.设在(
)
,a b 上
()
f
x
M
'
\le 
,则在(
)
,a b 上,
()
()
f x
f
y
M x
y
-
\le 
-
; 
特别地,当
()
f
x
'
在(
)
,a b 上有界时,
()
f x 在(
)
,a b 上有界. 
\textcolor{lyred}{例}. sin
sin
x
y
x
y
-
\le 
-
, cos
cos
x
y
x
y
-
\le 
-
, arctan
arctan
x
y
x
y
-
\le 
-
. 
\textcolor{lyred}{例}.设a
b
<
,则
(
)
(
)
a
b
a
b
e
b
a
e
e
e
b
a
-
<
-
<
-
. 
\textcolor{lyred}{例}.当
x >
时,
(
)
ln 1
x
x
x
x <
+
<
+
;特别地,
ln 1
n
n
n


<
+
<


+


. 
证.
(
)
(
)
ln 1
ln 1 0
x
x
x
\psi 
+
-
+
<
=
<
+
-
+
,证毕. 
\textcolor{lyred}{推论}.设在区间I 上
()
f
x
'
\equiv 
,则在I 上
()
f x 为常数. 
\textcolor{lyred}{推论}.设在(
)
,a b 内
()
()
f
x
g
x
'
'
\equiv 
,则
()
()
f x
g x
C
=
+
. 
\textcolor{lyred}{例}.设a
c
b
<
<
,
()
()
()
f a
f b
f c
=
<
,证明:
(
)
,a b
\psi 
\exists \in 
,使得
()
f
\psi 
''
<
. 
证.由中值定理,
(
)
,a c
\eta 
\exists \in 
,使得
()
()
()
f c
f a
f
c
a
\eta 
-
'
=
>
-
,同理,
(
)
,c b

\exists 
\in 
,使得 
()
()
()
f b
f c
f
b
c

-
'
=
<
-
,故
(
)
,
\psi 
\eta 
\exists \in 
,使得
()
()
()
f
f
f

\eta 
\psi 

\eta 
'
'
-
''
=
<
-
,证毕. 
\textcolor{lyred}{例}.设
()
()
f a
f b
=
=,证明:存在
(
)
,
,a b
\psi \eta \in 
,使得
()
()
f
f
e\psi \eta 
\eta 
\eta 
-
'
+
=
. 
证.
()
()
()
()
f
f
e
e
f
f
e
\psi \eta 
\eta 
\psi 
\eta 
\eta 
\eta 
\eta 
-
'
'
+
=
\Leftrightarrow 
+
=




,令
()
()
x
F x
e f x
=
,则 
(
)
,a b
\eta 
\exists 
\in 
,
(
)
,a b
\psi \in 
,使得
()
()
()
a
b
F a
F b
e
e
F
e
a
b
a
b
\psi 
\eta 
-
-
'
=
=
=
-
-
,证毕. 
\textcolor{lyred}{例}.
(
)
(
)
lim sin
sin
lim
cos
x
x
x
x
x
x
\psi 
\to +\infty 
\to +\infty 
+-
=
+-
\cdots 
=
. 
\textcolor{lyred}{例}.
lim
lim
2 ln 2
ln 2
x
x
x
x
x
x
x
x
\psi 
+
\to \infty 
\to \infty 




-
=
-
\cdots 
=




+




. 
\textcolor{lyred}{例}.
lim
arctan
arctan
lim
n
n
n
n
n
n
n
n
\psi 
\to \infty 
\to \infty 




-
=
-
\cdots 
=




+
+
+




. 
\textcolor{lyred}{例}.
(
)
(
)
(
)
(
)
(
)
(
)
tan tan
sin sin
tan tan
tan sin
tan sin
sin sin
lim
lim
lim
x
x
x
x
x
x
x
x
x
x
x
x
\to 
\to 
\to 
-
-
-
=
+
= 
tan
sin
tan
sin
limsec
lim
x
x
x
x
x
x
x
x
\psi 
\to 
\to 
-
-
\cdots 
+
=
+
=. 
\vspace{4pt}
{\bfseries \textcolor{lyred}{三.柯西中值定理}}

\textcolor{lyred}{定理}.设
()
f x , ()
g x 在[
]
,a b 上连续,(
)
,a b 内可导,且
()
g
x
'
\neq 
,则存在
(
)
,a b
\psi \in 
, 
使得
()
()
()
()
()
()
f
f b
f a
g
g b
g a
\psi 
\psi 
'
-
=
'
-
. 
\textcolor{lyred}{例}.设
a >
,证明:存在
(
)
,
,a b
\psi \eta \in 
,使得
()
()
abf
f
\psi 
\eta 
\eta 
'
'
=
. 
证.
(
)
,a b
\psi 
\exists \in 
,
()
()
()
f a
f b
f
a
b
\psi 
-
'
=
-
,上式
()
()
()
f a
f b
f
a
b
\eta 
\eta 
'
-
\Leftrightarrow 
=
-
-
,令()
g x
x
=
, 
则
(
)
,a b
\eta 
\exists 
\in 
,使得
()
()
()
()
()
()
f a
f b
f
g a
g b
g
\eta 
\eta 
'
-
=
'
-
,即得,证毕. 
\textcolor{lyred}{例}.设
a b
\cdots 
>
,证明:存在
(
)
,
,a b
\psi \eta \in 
,使得
()
(
)
()
a
b f
f
\eta 
\psi 
\eta 
'
+
'
=
. 
证.
(
)
,a b
\psi 
\exists \in 
,
()
()
()
f a
f b
f
a
b
\psi 
-
'
=
-
,上式
()
()
()
f a
f b
f
a
b
\eta 
\eta 
'
-
\Leftrightarrow 
=
-
,令()
g x
x
=
, 
则
(
)
,a b
\eta 
\exists 
\in 
,使得
()
()
()
()
()
()
f a
f b
f
g a
g b
g
\eta 
\eta 
'
-
=
'
-
,即得,证毕. 
\textcolor{lyred}{补充练习 }
1.设
(
)
()
f
f
-
=
=
,
()
f
=
,证明:曲线
()
y
f x
=
有一条切线,它平行于直线 
x
y
-
+
=
. 
证.令
()
()
x
F x
f x
=
-
,则
(
)
F -
=,
()
F
=
,
()
F
=-,故
()
F x 有最大值点 
(
)
2,2
\psi \in -
,于是
()
F \psi 
'
=
,证毕. 
2.证明:方程
x
x
-
+=
在(
)
0,1 内有且仅有一个根. 
证.设
()
f x
x
x
=
-
+,则
()
f
=,
()1
f
=-,故在(
)
0,1 内有零点; 
又在(
)
0,1 上
()
f
x
x
'
=
-
\neq 
,故至多只有一个零点,证毕. 
3.设
()
f
x
'
\neq ,
()
f x
<
<,证明:存在唯一的
(
)
0,1
\psi \in 
,使得
()
f \psi 
\psi 
=
. 
证.令
()
()
F x
f x
x
=
-
,则
()
F
>
,
()1
F
<
,故存在
(
)
0,1
\psi \in 
,使得
()
F \psi =
; 
又在(
)
0,1 内
()
()1
F
x
f
x
'
'
=
-\neq 
,故\psi 唯一,证毕. 
4.证明:存在
(
)
0,
\psi 
\pi 
\in 
,使得
()
()
cot
f
f
\psi 
\psi 
\psi 
'
\cdots 
+
=
. 
证.
()
()
()
cos
lnsin
ln
sin
f
x
x
x
f x
x
f x
'
'
+
=
\Leftrightarrow 
+
=




()
sin
x f x
'
\Leftrightarrow 
\cdots 
=




; 
令
()
()
sin
F x
x f x
=
\cdots 
,则
()
()
F
F \pi 
=
=
,证毕. 
5.设
()
f x 有界,且
()
lim
x
f
x
A
\to +\infty 
'
=
,证明:
A =
. 
证.
(
)
(
)
()
n
f
n
f n
f
n
\psi 
-
'
=
,其中
(
)
,2
n
n
n
\psi \in 
,得
(
)
lim
n
n
f
\psi 
\to \infty 
'
=
,由于lim
n
n
\psi 
\to \infty 
=+\infty , 
故
(
)
()
lim
lim
n
n
x
f
f
x
A
\psi 
\to \infty 
\to +\infty 
'
'
=
=
,于是
A =
.证毕. 
6.设
()
f
x
'
\neq 
,证明:存在
(
)
,
,a b
\psi \eta \in 
,使得
()
()
a
b
f
e
e e
f
a
b
\eta 
\psi 
\eta 
-
'
-
=
'
-
. 
证.
(
)
,a b
\psi 
\exists \in 
,
()
()
()
f a
f b
f
a
b
\psi 
-
'
=
-
,上式
()
()
()
a
b
f a
f b
f
e
e
e\eta 
\eta 
'
-
\Leftrightarrow 
=
-
,令()
x
g x
e
=
, 
则
(
)
,a b
\eta 
\exists 
\in 
,
()
()
()
()
()
()
f a
f b
f
g a
g b
g
\eta 
\eta 
'
-
=
'
-
,即得,证毕. 
\vspace{6pt}
{\large \bfseries \textcolor{lyred}{第3.2 节 洛必达法则}}

\vspace{4pt}
{\bfseries \textcolor{lyred}{一.无穷小或无穷大的商}}

\textcolor{lyred}{定理}.设(1)
()
()
lim
lim
x
a
x
a
f x
g x
\to 
\to 
=
=
;(2)在某个
()
U a
\circ 
内,
()
f x 与()
g x 均可导,且 
()
g
x
'
\neq 
;(3)
()
()
lim
x
a
f
x
g
x
\to 
'
'
存在,或为\infty ,则
()
()
()
()
lim
lim
x
a
x
a
f x
f
x
g x
g
x
\to 
\to 
'
=
'
. 
\textcolor{lyred}{定理}.设(1)
()
()
lim
lim
x
x
f x
g x
\to \infty 
\to \infty 
=
=
;(2)当x
X
>
时,
()
f x 与()
g x 均可导,且 
()
g
x
'
\neq 
,(3)
()
()
lim
x
f
x
g
x
\to \infty 
'
'
存在,或为\infty ,则
()
()
()
()
lim
lim
x
x
f x
f
x
g x
g
x
\to \infty 
\to \infty 
'
=
'
. 
\textcolor{lyred}{注}.对于\infty 
\infty 型未定式,若
()
()
lim
x
f
x
g
x
\to 
'
'
\Omega 
存在,或为\infty ,则
()
()
()
()
lim
lim
x
x
f x
f
x
g x
g
x
\to 
\to 
'
=
'
\Omega 
\Omega 
. 
\textcolor{lyred}{例}.
lim
lim
lim
x
x
x
x
x
x
x
x
x
x
x
x
x
\to 
\to 
\to 
-
+
-
=
=
=
-
-
+
-
-
-
. 
\textcolor{lyred}{例}.
(
)
1 sin
cos
lim sec
tan
lim
lim
cos
sin
x
x
x
x
x
x
x
x
x
\pi 
\pi 
\pi 
\to 
\to 
\to 
-
-
-
=
=
=
-
. 
\textcolor{lyred}{例}.
tan
tan
1 sec
2sec
tan
lim
lim
lim
lim
sin
x
x
x
x
x
x
x
x
x
x
x
x
x
x
x
x
\to 
\to 
\to 
\to 
-
-
-
-
=
=
=
=-
. 
\textcolor{lyred}{例}.
arcsin
sin
sin
limcot
lim
lim
lim
arcsin
sin
x
x
x
x
x
x
u
u
u
u
x
x
x
x
u
u
\to 
\to 
\to 
\to 
-
-
-


-
=
=
=
=-




. 
\textcolor{lyred}{例}.
(
)
lim
lim
lim
1 sin
tan
x
x
x
x
x
x
x
x
x
e
e
e
e
e
e
x
x
x
x
-
-
-
\to 
\to 
\to 
-
-
+
=
=
=\infty 
+
. 
\textcolor{lyred}{例}.
(
)(
)
sin
cos
sin
cos
sin
cos
lim
lim
sin
x
x
x
x
x
x
x
x
x
x
x
x
x
x
\to 
\to 
+
-
-
=
= 
sin
cos
sin
cos
sin
cos
sin
lim
lim
2lim
2lim 3
x
x
x
x
x
x
x
x
x
x
x
x
x
x
x
x
x
x
x
\to 
\to 
\to 
\to 
+
-
-
\cdots 
=
=
=
. 
\textcolor{lyred}{例}.
(
)
(
)
sin
sinsin
sin
cos
cos sin
cos
sin
sinsin
lim
lim
lim
x
x
x
x
x
x
x
x
x
x
x
x
x
x
\to 
\to 
\to 
-
-
\cdots 
-
=
=
= 
(
)
1 cos sin
sin
lim
lim
x
x
x
x
x
x
\to 
\to 
-
=
=
;或者,直接原式
sin
lim
x
x
x
x
\to 
-
=
=
. 
\textcolor{lyred}{例}.
(
)
ln
lim
lim
lim
x
x
x
x
x
x
x
x
\alpha 
\alpha 
\alpha 
\alpha 
\alpha 
\alpha 
-
-
\to +\infty 
\to +\infty 
\to +\infty 
>
=
=
=
,
ln
lim
lim
n
n
n
n
n
n
e
\to \infty 
\to \infty 
=
=. 
\textcolor{lyred}{例}.
(
)
(
)
lim
0,
lim
lim
x
x
x
x
x
x
x
x
x
e
e
e
\alpha 
\alpha 
\alpha 
\lambda 
\lambda 
\lambda 
\alpha \alpha 
\alpha 
\alpha 
\lambda 
\lambda 
\lambda 
-
-
\to +\infty 
\to +\infty 
\to +\infty 
-
>
>
=
=
=
=
\Lambda 
. 
\textcolor{lyred}{例}.设
(
)
f
x
''
存在,求
(
)
(
)
(
)
lim
h
f x
h
f x
h
f x
h
\to 
+
+
-
-
. 
解.
(
)
f
x
''
存在
()
f
x
'
\Rightarrow 
在
0x 的某个邻域内存在,故 
(
)
(
)
(
)
(
)
(
)
lim
lim
h
h
f x
h
f x
h
f x
f
x
h
f
x
h
h
h
\to 
\to 
'
'
+
+
-
-
+
-
-
=
= 
(
)
(
)
(
)
(
)
(
)
lim
lim
h
h
f
x
h
f
x
f
x
h
f
x
f
x
h
h
\to 
\to 
'
'
'
'
+
-
-
-
''
-
=
. 
\textcolor{lyred}{注}.最后一步不能用洛必达法则,除非
()
f
x
''
存在,且在
x
x
=
处连续. 
\textcolor{lyred}{例}.设
()
f
=
,
()
()
f
f
'
''
=
=,令()
(),
1,
f x
x
g x
x
x

\neq 

=

=

,求
()
g'
. 
解.
()
()
()
()
()
()
lim
lim
lim
x
x
x
g x
g
f x
x
f
x
g
f
x
x
x
\to 
\to 
\to 
'
-
-
-
'
''
=
=
=
=
=
-
. 
\textcolor{lyred}{注}.最后一步不能用洛必达法则,除非
()
f
x
''
存在,且在
x =
处连续. 
\textcolor{lyred}{例}.
sin
1 cos
lim
lim
cos
1 sin
x
x
x
x
x
x
x
x
\to \infty 
\to \infty 
+
+
\Rightarrow 
-
+
,不存在(非常数的周期函数),而
sin
lim
cos
x
x
x
x
x
\to \infty 
+
=
-
. 
\textcolor{lyred}{例}.
sin
2 sin
cos
lim
lim
x
x
x
x
x
x
x
x
\to 
\to 
-
\Rightarrow 
,不存在,而
sin
lim
x
x
x
x
\to 
=
. 
\vspace{4pt}
{\bfseries \textcolor{lyred}{二.其它未定式(0,,}}

\textcolor{lyred}{0}\textcolor{lyred}{ }\textcolor{lyred}{,}\textcolor{lyred}{1}\textcolor{lyred}{\infty }\textcolor{lyred}{,}
\textcolor{lyred}{\infty }\textcolor{lyred}{) }
\textcolor{lyred}{例}.
arctan
lim
arctan
lim
lim
lim
x
x
x
x
x
x
x
x
x
x
x
x
\pi 
\pi 
-
-
\to +\infty 
\to +\infty 
\to +\infty 
\to +\infty 
-
-


+
-
=
=
=
=


-
+


. 
\textcolor{lyred}{例}.
(
)
ln
lim
ln
lim
lim
lim
x
x
x
x
x
x
x
x
x
x
x
\alpha 
\alpha 
\alpha 
\alpha 
\alpha 
\alpha 
\alpha 
+
+
+
+
-
-
-
-
\to 
\to 
\to 
\to 
>
=
=
=-
=
-
; 
或者,
(
)
ln
lim
ln
lim
t
x
t
x
t
x
x
t
\alpha 
\alpha 
\alpha 
+
=
\to +\infty 
\to 
-
>
=
=
(倒代换). 
\textcolor{lyred}{例}.
lim
ln
lim
x
x
x
x
x
x
e
e
+
\to 
+
\to 
=
=
=,
(
)
(
)
ln
ln
lim 1
lim
x
x
x
x
x
x
x
x
e
+
+
\to 
\to 


+
=
+
=
=




. 
\textcolor{lyred}{例}.
lim
lim
t
n
x
n
t
x
e
x e
t
+
\to +\infty 
\to 
=
=\infty ,
lim
lim
n
x
n
t
t
x
e
t
x
e
+
-
\to +\infty 
\to 
=
=
. 
\textcolor{lyred}{例}.
(
)
ln
ln
1 1
ln
lim
lim
lim
lim
ln
1 ln
ln
ln
x
x
x
x
x
x
x
x
x
x
x
x
x
x
x
x
x
x
x
x
x
\to 
\to 
\to 
\to 
-
+
+-


-
=
=
=
=


-
-
-
+
-


+
. 
\textcolor{lyred}{例}.
(
)
ln
lim
ln
lim
ln
lim
x
t
t
t
t
x
x
t
t
t
+
+
\to +\infty 
\to 
\to 
+


-
=
-
=
=+\infty 




. 
\textcolor{lyred}{例}.
(
)
(
)
ln 1
lim
ln 1
lim
ln 1
lim
t
x
x
t
t
t
t
x
x
t
x
t
t
t
+
+
=
\to +\infty 
\to 
\to 
-
+






-
+
=
-
+
=
=












. 
\textcolor{lyred}{例}.
(
)
(
)
(
)
(
)
(
)
(
)
lim ln ln 1
lim ln 1
1 ln 1
lim 1
ln 1
ln
lim 1
x
x
x
x
x
x
x
x
x
x
x
x
e
e
e
e
-
-
-
\to 
\to 
\to 
-
-
+-
-
-
-
-
\to 
-
=
=
=
=
=. 
\textcolor{lyred}{例}.
ln
arctan
lim
arctan
lim
arctan
lim
x
x
x
x
x
x
x
x
x
e
e
e
\pi 
\pi 
\pi 
\pi 
\to +\infty 




\cdots 
-
-








\to +\infty 
\to +\infty 

=
=
=




. 
\textcolor{lyred}{例}.
tan
tan
tan
lim
lim
ln
tan
lim
tan
lim
x
x
n
x
x
x x
x
x
x
x
x
x
n
x
x
n
e
e
e
e
n
x
+
+
\to 
\to 
+
-


-




\to \infty 
\to 




\cdots 
=
=
=
=
=








. 
\textcolor{lyred}{例}.
(
)
ln 1
lim
ln 1
lim
lim
lim 1
t
t
x
t
x
t
t
x
x
x
t
x
e
e
e
e
x
+
\to +\infty 
\to +\infty 
\to 
+
=
+


+
+




\to 


+
=
=
=
=
=




. 
\textcolor{lyred}{补充练习 }
1.
tan
sec
cos 3
3sin 6
6cos6
lim
lim
lim
lim
lim
tan3
3sec 3
cos
sin 2
2cos2
x
x
x
x
x
x
x
x
x
x
x
x
x
x
x
\pi 
\pi 
\pi 
\pi 
\pi 
\to 
\to 
\to 
\to 
\to 
=
=
=
=
=. 
2.
1 sin 2
cos
sin
cos
lim
lim
lim
sin
sin
x
x
x
x
x
x
x
x
x
x
x
x
x
x
\to 
\to 
\to 
-


-
-
=
=
=




sin 4
2cos4
8sin 4
lim
lim
lim
x
x
x
x
x
x
x
x
x
x
\to 
\to 
\to 
-
-
=
=
=
. 
3.
(
)
(
)
ln 1
lim
lim
lim
x
n
x
x
n
x
x
x
e
e
n
e
e
e
n
x
x
+
-
\to \infty 
\to 
\to 


+
-
-


+
-
=
=
=-










. 
4.
(
)
sin
sin
sin
sin
lim
lim
lim
lim
x
x
x
x
x
x
x
x
x
x
x
x
x
x
x
x
x
x
x
x
+
+
+
+
\to 
\to 
\to 
\to 




-
-




-
-




=
=
=
=-
. 
5.
(
)
(
)
1 ln
lim
lim
lim
lim
ln
ln
ln
ln
x
x
x
x
x
x
x
x
x
x
x
x
e
x
x
x
x
x
x
x
x
-
-
\to 
\to 
\to 
\to 
-
-
-
-
=
=
=
=-
-
+
-
+
-
+
-
+
. 
6.
(
)
2 2cos
2 2cos
2 2cos
2cos
lim
lim
lim
x
x
x
x
x
x
x
x
e
e
e
e
x
x
x
x
x
-
-+
-
\to 
\to 
\to 
-
-
-+
=
=
=
. 
7.设
lim
x
x
x
ax
b
x
\to \infty 




+
-
-
=










,求a ,b . 
解.
lim 1
x
x
a
e
x
\to \infty 


=
+
=




,
(
)
lim
lim
x
t
x
t
t
e
b
x
ex
x
t
\to \infty 
\to 


+
-


+
=
+
-
=
=










(
)
(
)
(
)
ln 1
ln 1
ln 1
lim
lim
lim
t
t
t
t
t
t
t
t
t
e
e
e
e
e
e
t
t
t
+
+
+
-
\to 
\to 
\to 
+
-
-
-
=
=
=-
,故
e
b =-
-
. 
8.设
()
sin 6
lim
x
xf
x
x
x
\to 
+
=
,求
()
lim
x
f x
x
\to 
+
. 
解.
()
()
sin 6
sin 6
sin 6
lim
lim
lim
x
x
x
f x
xf x
x
x
x
x
x
x
x
x
\to 
\to 
\to 
+
+
-
+
-
=
=
-
=
. 
\vspace{6pt}
{\large \bfseries \textcolor{lyred}{第3.3 节 泰勒公式}}

\vspace{4pt}
{\bfseries \textcolor{lyred}{一.泰勒公式}}

设
(
)
( )
n
f
x
存在,记 
()
(
)
(
)(
)
(
)(
)
(
)(
)
( )
2!
!
n
n
n
f
x
f
x
P
x
f
x
f
x
x
x
x
x
x
x
n
''
'
=
+
-
+
-
+
+
-
\Lambda 
, 
称为\textcolor{lyred}{n 阶泰勒多项式},记
()
()
()
n
n
R
x
f x
P x
=
-
. 
\textcolor{lyred}{定理}.设
()(
)
n
f
x
存在,则当
x
x
\to 
时,
()
(
)
n
n
R
x
o
x
x


=
-

. 
\textcolor{lyred}{注}.
()
()
(
)
n
n
f x
P x
o
x
x


=
+
-

,称为
()
f x 的带\textcolor{lyred}{Peano 型余项的n 阶泰勒公式}. 
\textcolor{lyred}{定理(泰勒中值定理)}.设
()
f x 在含
0x 的(
)
,a b 内
n +阶可导,则
(
)
,
x
a b
\forall \in 
,存在 
介于
0x 与x 之间的\psi ,使得
()
(
)()
(
)(
)
1 !
n
n
n
f
R
x
x
x
n
\psi 
+
+
=
-
+
. 
\textcolor{lyred}{注}.
()
()
(
)()
(
)(
)
1 !
n
n
n
f
f
x
P
x
x
x
n
\psi 
+
+
=
+
-
+
称为
()
f x 的带\textcolor{lyred}{Lagrange 型余项的n 阶} 
\textcolor{lyred}{泰勒公式}. 
\vspace{4pt}
{\bfseries \textcolor{lyred}{二.麦克劳林公式}}

()
()
()
()
()()
(
)(
)
(
)
2!
!
1 !
n
n
n
n
f
f
f
x
f x
f
f
x
x
x
x
n
n
\theta 
+
+
''
'
=
+
+
+
+
+
+
\Lambda 
,其中0
\theta 
<
<; 
()
()
()
()
()()
(
)
2!
!
n
n
n
f
f
f x
f
f
x
x
x
o x
n
''
'
=
+
+
+
+
+
\Lambda 
. 
\textcolor{lyred}{例}.设
()
x
f x
e
=
,则
()()
n
f
=,故
()
2!
!
n
x
n
x
x
e
x
R
x
n
=+
+
+
+
+
\Lambda 
,其中 
()
(
)
(
)
1 0
1 !
x
n
n
e
R
x
x
n
\theta 
\theta 
+
=
<
<
+
,特别地,
1 1
2!
!
n
e
R
n
=++
+
+
+
\Lambda 
,其中 
(
)1 !
n
e
R
n
<
+
,于是,
n
\forall >,
!
e n
\cdots 
不可能是整数,即e 为无理数. 
\textcolor{lyred}{例}.设
()
sin
f x
x
=
,则
()()
sin 0
sin
n
n
n
f
\pi 
\pi 


=
+
=




,故 
(
)
(
)
()
sin
3!
5!
7!
1 !
m
m
m
x
x
x
x
x
x
R
x
m
-
-
=
-
+
-
+
+-
+
-
\Lambda 
,其中 
()
(
)
(
)
(
)
(
)
sin
1 2
1 !
m
m
m
x
m
R
x
x
o x
m
\pi 
\theta 
\theta 
+


+
+




=
<
<
=
+
. 
\textcolor{lyred}{注}.类似地,
(
)(
)
(
)
cos
2!
4!
!
m
m
m
x
x
x
x
o x
m
+
=-
+
-
+-
+
\Lambda 
; 
(
)
(
)
(
)
ln 1
n
n
n
x
x
x
x
x
o x
n
-
+
=
-
+
-
+-
+
\Lambda 
; 
(
)
(
)
(
)
(
)
(
)
2!
!
n
n
n
x
x
x
x
o x
n
\alpha 
\alpha \alpha 
\alpha \alpha 
\alpha 
\alpha 
-
-
-
+
+
=+
+
+
+
+
\Lambda 
\Lambda 
. 
\vspace{4pt}
{\bfseries \textcolor{lyred}{三.泰勒公式的应用}}

\textcolor{lyred}{例}.设
()
lim
x
f x
x
\to 
=,且
()
f
x
''
>
,证明:当
x \neq 
时,
()
f x
x
>
. 
证.
()
f
=
,
()
f '
=,故
()
()
()
()
2!
f
f x
f
f
x
x
x
\psi 
''
'
=
+
+
>
,证毕. 
\textcolor{lyred}{例(Jensen 不等式)}.
()
()
(
)
n
n
f x
f x
x
x
f
x
f
n
n
+
+
+
+


''
\ge 
\Rightarrow 
\ge 




\Lambda 
\Lambda 
. 
证.令
n
x
x
x
n
+
+
=
\Lambda 
,则
()
(
)
(
)(
)
()(
)
i
i
i
i
f
f x
f x
f
x
x
x
x
x
\psi 
''
'
=
+
-
+
-
,故 
()
(
)
(
)(
)
()
(
)
n
i
i
i
i
f x
f x
f
x
x
x
f x
nf x
=
'
\ge 
+
-
\Rightarrow 
\ge 
\sum 
,即得,证毕. 
\textcolor{lyred}{例}.
(
)
(
)
3!
2!
sin
cos
lim
lim
sin
x
x
x
x
x
o x
x
o x
x
x
x
x
x
\to 
\to 


-
+
-
-
+


-


=
=-
+
=
. 
\textcolor{lyred}{例}.
(
)
(
)
(
)
(
)
2!
4!
2!
cos
lim
lim
ln 1
x
x
x
x
x
x
x
o x
o x
x
e
x
x
x
x
o x
x
-
\to 
\to 




-
+
+
---
-
-
-




-




=
=
+
-
-
+
-
. 
\textcolor{lyred}{例}.设
()
sin 6
lim
x
xf x
x
x
\to 
+
=
,求
()
lim
x
f
x
x
\to 
+
. 
解.
()
(
)
(
)
()
()
3!
lim
lim
lim
x
x
x
x
xf x
x
o x
f x
f x
x
x
x
\to 
\to 
\to 


+
-
+


+
+



=
-
\Rightarrow 
=
. 
\textcolor{lyred}{例}.设
(
)
ln 1
lim
x
x
ax
bx
x
\to 
+
-
+
=
,求a ,b . 
解.
(
)
lim
x
x
x
o x
ax
bx
a
x
\to 
-
+
-
+
=
\Rightarrow 
=,
b =
. 
\textcolor{lyred}{例}.设
()
f ''
存在,
()
(
)
ln 1
lim
x
xf x
x
x
\to 
-
+
=
,求
()
f
,
()
f '
,
()
f ''
. 
解.
()
()
()
(
)
(
)
lim
x
f
x
x
x f
f
x
x
o x
x
o x
x
\to 
''




'
+
+
+
-
-
+
+








=
,故 
()
f
=,
()
f '
=-
,
()
f ''
=
. 
\vspace{6pt}
{\large \bfseries \textcolor{lyred}{第3.4 节 函数的单调性与曲线的凹凸性}}

\vspace{4pt}
{\bfseries \textcolor{lyred}{一.函数的单调性}}

{\bfseries \textcolor{lyred}{1.单调性的判定}}

\textcolor{lyred}{定理}.设
()
f x 在[
]
,a b 上连续,在(
)
,a b 上可导,则 
(1)当在(
)
,a b 上
()
f
x
'
>
时,
()
f x 在[
]
,a b 上单调增加; 
(2)当在(
)
,a b 上
()
f
x
'
<
时,
()
f x 在[
]
,a b 上单调减少. 
\textcolor{lyred}{推论}.设
()
f x 在[
]
,a b 上连续,(1)若它在(
)
,a b 上除了有限个点外满足
()
f
x
'
>
, 
则它在[
]
,a b 上单调增加;(2)若它在(
)
,a b 上除了有限个点外满足
()
f
x
'
<
,则它 
在[
]
,a b 上单调减少. 
\textcolor{lyred}{定理}.设
()
f x 在[
]
,a b 上连续,(
)
,a b 上可导,且在[
]
,a b 的任何子区间上不为常数, 
则(1)当在(
)
,a b 上
()
f
x
'
\ge 
时,
()
f x 单调增加; 
(2)当在(
)
,a b 上
()
f
x
'
\le 
时,
()
f x 单调减少. 
\textcolor{lyred}{例}.
sin
y
x
x
=
-
,
cos
y
x
'=-
\ge 
,故
sin
y
x
x
=
-
单增. 
\textcolor{lyred}{例}.讨论
y
x
=
的单调性. 
解.
y
x
-
'=
,故在(
]
,0
-\infty 
上单减,[
)
0,+\infty 上单增,
x =
为不可导点. 
\textcolor{lyred}{例}.讨论
x
y
e
x
=
-
-的单调性. 
解.
x
y
e
'=
-,故在(
]
,0
-\infty 
上单减,在[
)
0,+\infty 上单增,
x =
为驻点. 
\textcolor{lyred}{注}.函数单调区间的分界点总是驻点或不可导点. 
\textcolor{lyred}{例}.确定
()
f x
x
x
x
=
-
+
-的单调区间. 
解.
()
(
)(
)
f
x
x
x
x
x
x
'
=
-
+
=
-
-
=
\Rightarrow 
=,2 ,列表如下 
x       (
)
,1
-\infty 
    1     (
)
1,2    2     (
)
2,+\infty  
()
f
x
'
     +      0      -      0       + 
故
()
f x 在(
]
,1
-\infty 
和[
)
2,+\infty 上单调增加,在[
]
1,2 上单调减少. 
\textcolor{lyred}{例}.确定
()
ln
f x
x
x
=
-
的单调区间. 
解.
()
(
)
f
x
x
x
x
x
'
=
-
=
-
=
\Rightarrow 
=,加上间断点
x =
,列表如下: 
x       (
)
,0
-\infty 
   0     (
)
0,1    1     (
)
1,+\infty  
()
f
x
'
     +      \times       -      0      + 
故
()
f x 在(
)
,0
-\infty 
上单调增加,在(
]
0,1 上单调减少,在[
)
1,+\infty 上单调增加. 
{\bfseries \textcolor{lyred}{2.不等式的证明}}

\textcolor{lyred}{例}.证明:当
x >时,
x
x
>
-
. 
证.令
()
f x
x
x


=
-
-




,则
()
f
x
x
x
'
=
-
,由于当
x >时,
x
x
x
<
<
, 
x
x
>
,得
()
f
x
'
>
,故
()
f x 在[
)
1,+\infty 上单增,而
()1
f
=
,证毕. 
\textcolor{lyred}{例}.证明:当0
x
\pi 
<
<
时,sin
tan
x
x
x
+
>
. 
证.
()
sin
tan
f x
x
x
x
=
+
-
,
()
cos
sec
cos
cos
f
x
x
x
x
x
'
=
+
-
>
+
-
= 
cos
cos
x
x


-
>




,故
()
f x 在0, 2
\pi 





上单增,而
()
f
=
,证毕. 
\textcolor{lyred}{例}.证明:当
x >
时,
2x
x
>
. 
证.
ln2
2ln
x
x
x
x
>
\Leftrightarrow 
>
,令
()
ln 2
2ln
f x
x
x
=
-
,则
()
ln 2
f
x
x
'
=
-
,当
x >
时, 
()
ln 2
ln
f
x
e
'
>
-
=
>
,故
()
f x 在[
)
4,+\infty 上单增,而
()
f
=
,证毕. 
\textcolor{lyred}{例}.证明:当0
x
y
\pi 
<
<
<
时, tan
tan
y
y
x
x
>
. 
证.只需证tan
tan
y
x
y
x
>
,即tan x
x
单增;
()
tan
sec
tan
g x
x
x
x
x
x
x
x
'
-

=
=




, 
在0, 2
\pi 






上,
()
()
2 sec
tan
g
x
x
x
x
g x
'
=
>
\Rightarrow 
在0, 2
\pi 





上单增,而()
g
=
,故 
在0, 2
\pi 






上()
g x >
,即得,证毕. 
\textcolor{lyred}{例}.设
()
f x 二阶可导,
()
f
=
,
()
f
x
''
>
,证明:
()
f x
x
在(
)
0,+\infty 上单调增加. 
证.设
()
()
f x
F x
x
=
,则
()
()
()
()
xf
x
f x
g x
F
x
x
x
'
-
'
=
=
,在(
)
0,+\infty 上, 
()
()
()
g
x
xf
x
g x
'
''
=
>
\Rightarrow 
在[
)
0,\infty 上单增,而()
g
=
,故在(
)
0,+\infty 上()
g x >
, 
即得,证毕. 
\textcolor{lyred}{例}.证明:当
x >
时,
x
e
e
x
\ge 
,当且仅当x
e
=
时,等号成立. 
证.
ln
x
e
e
x
x
e
x
\ge 
\Leftrightarrow 
\ge 
,令()
ln
f x
x e
x
=-
,则
()
e
f
x
x
'
=-
,在0
x
e
<
<
上
()
f
x
'
<
, 
在x
e
>
上
()
f
x
'
>
,故
()
f e =
为最小值,证毕. 
{\bfseries \textcolor{lyred}{3.方程根的个数}}

\textcolor{lyred}{例}.求方程
xe
x
-
+
=
的实根个数. 
解.设
()
2,
2,
x
x
x
e
x
x
f x
e
x
e
x
x

+
+
\le -
=
-
+
=
-
-
>-

,则
()
1,
1,
x
x
e
x
f
x
e
x

+
<-
'
=
-
>-

,于是 
()
f x 在(
]
, 2
-\infty -
上单增,在[
]
2,0
-
上单减,在[
)
0,+\infty 上单增,又 
()
lim
x
f x
\to -\infty 
=-\infty ,
(
)
f
e
-
=
>
,
()
f
=-<
,
()
lim
x
f x
\to +\infty 
=+\infty ,故
()
f x 有三个 
零点,即
xe
x
-
+
=
有三个实根. 
\textcolor{lyred}{例}.确定曲线
ln
y
x
x
=
+
与
4ln
y
x
k
=
+
的交点个数. 
解.令
()(
)(
)
ln
4ln
f x
x
x
x
k
=
+
-
+
,则
()
(
)
4 ln
x
x
f
x
x
+
-
'
=
,
()1
f '
=
, 
当0
x
<
<时,
()
f
x
'
<
,当
x >时,
()
f
x
'
>
,故
()1
f
k
=
-
为最小值,并且 
()
lim
x
f x
\to +\infty 
=+\infty ,
()
lim
x
f x
+
\to 
=+\infty ,(1)当
k <
时,
()
f x 无零点,故曲线无交点; 
(2)当
k =
时,
()
f x 有唯一零点,故曲线有一个交点; 
(3)当
k >
时,
()
f x 有两个零点,故曲线有两个交点. 
\vspace{4pt}
{\bfseries \textcolor{lyred}{二.曲线的凹凸性与拐点}}

\textcolor{lyred}{定义}.设
()
f x 在区间I 上连续,若
x
x
I
\forall 
\neq 
\in 
,恒有
(
)
(
)
f x
f x
x
x
f
+
+

<




, 
则称曲线
()
y
f x
=
为\textcolor{lyred}{凹弧}; 
若
x
x
I
\forall 
\neq 
\in 
,恒有
(
)
(
)
f x
f x
x
x
f
+
+

>




,则称曲线
()
y
f x
=
为\textcolor{lyred}{凸弧}. 
\textcolor{lyred}{定理}.设
()
f x 在[
]
,a b 上连续,在(
)
,a b 上可导,则 
(1)当在(
)
,a b 上
()
f
x
'
单调增加时,曲线
()
y
f x
=
是凹弧; 
(2)当在(
)
,a b 上
()
f
x
'
单调减少时,曲线
()
y
f x
=
是凸弧. 
\textcolor{lyred}{定理}.设
()
f x 在[
]
,a b 上连续,在(
)
,a b 上二阶可导,则 
(1)当在(
)
,a b 上
()
f
x
''
>
时,曲线
()
y
f x
=
是凹弧; 
(2)当在(
)
,a b 上
()
f
x
''
<
时,曲线
()
y
f x
=
是凸弧. 
\textcolor{lyred}{例}.判断曲线
ln
y
x
=
的凹凸性. 
解.
y
x-
''=-
<
,故曲线
ln
y
x
=
在(
)
0,+\infty 内是凸的. 
\textcolor{lyred}{例}.判断曲线
x
y =
的凹凸性. 
解.
y
x
'=
,
y
x
''=
,曲线在(
]
,0
-\infty 
上为凸,[
)
0,+\infty 上为凹,(
)
0,0 为拐点. 
\textcolor{lyred}{例}.判断曲线
y
x
=
的凹凸性. 
解.
y
x
-
''=-
,曲线在(
]
,0
-\infty 
上为凹,在[
)
0,+\infty 上凸,(
)
0,0 为拐点. 
\textcolor{lyred}{定义}.若连续曲线
()
y
f x
=
在点(
)
,
x
y
两侧有不同的凹凸性,则称该点为曲线的 
\textcolor{lyred}{拐点}. 
\textcolor{lyred}{注}.在拐点(
)
,
x
y
处总有
(
)
f
x
''
=
,或
(
)
f
x
''
不存在;反之不对. 
\textcolor{lyred}{例}.求曲线
y
x
x
=
-
+的凹凸区间及拐点. 
解.
y
x
x
'=
-
,
y
x
x
x x


''=
-
=
-




,
y
x
''=
\Rightarrow 
=
, 2
3 ,列表如下: 
x       (
)
,0
-\infty 
0, 3






3      2 ,


+\infty 




y''        +        0       -        0         + 
图形       凹      拐点      凸       拐点       凹 
故在(
]
,0
-\infty 
, 2 ,


+\infty 


上凹,在
0, 3






上凸,(
)
0,1 和2 11
,
3 27






为拐点. 
\textcolor{lyred}{例}.求曲线
(
)
x
y
x
=
-
的凹凸区间及拐点. 
解.
(
)
(
)
x
x
y
x
-
'=
-
,
(
)
x
y
x
''=
-
,
y
x
''=
\Rightarrow 
=
,又
()
y''
不存在,列表如下: 
x       (
)
,0
-\infty 
       0      (
)
0,1       1      (
)
1,+\infty  
y''        -          0       +        \times         + 
图形       凸        拐点      凹        \times         凹 
故凸区间为(
]
,0
-\infty 
,凹区间为[
)
0,1 和(
)
1,+\infty ,拐点为(
)
0,0 . 
\textcolor{lyred}{定理}.设
(
)
f
x
''
=
,
(
)
f
x
'''
\neq 
,则
(
)
(
)
,
x
f x
是
()
y
f x
=
的拐点. 
\textcolor{lyred}{例}.证明:(
)
(
)
ln
ln
ln
0,
x
y
x
y
x
x
y
y x
y
+
+
\le 
+
>
>
. 
证.只要证明
ln
ln
ln
x
y
x
y
x
x
y
y
+
+
+
\le 
,令
()
(
)
ln
f x
x
x x
=
>
,则 
()
ln
f
x
x
'
=
+,
()
f
x
x
''
=
>
,故
()
y
f x
=
是凹弧,即得,证毕. 
\textcolor{lyred}{补充练习 }
1.证明:当0
x
\pi 
<
<
时,
sin
x
x
x
>
-
. 
证.令
()
sin
f x
x
x
x
=
-
+
,则
()
cos
f
x
x
x
'
=
-+
,
()
sin
f
x
x
x
''
=-
+
, 
在(
)
0,2\pi 上
()
()
f
x
f
x
''
'
>
\Rightarrow 
在[
]
0,2\pi 上单增,
()
f '
=
,故在(
)
0,2\pi 上, 
()
()
f
x
f x
'
>
\Rightarrow 
在[
]
0,2\pi 上单增,而
()
f
=
,即得,证毕. 
2.设,
p q >, 1
p
q
+
=,证明:当
x >
时,
p
x
x
p
q
+
\ge 
. 
证.令
()
p
x
f x
x
p
q
=
+
-
,则
()
p
f
x
x
x
-
'
=
-=
\Rightarrow 
=,又当0
x
<
<时, 
()
f
x
'
<
,当
x >时,
()
f
x
'
>
,故
()1
f
=
为最小值,证毕. 
3.设
a
b
>
>
,
()
f x
x
ax
b
=
-
+
,求方程
()
f x =
的实根个数. 
解.
()
f
x
x
a
x
a
'
=
-
=
\Rightarrow 
=\pm 
,故
()
f x 在(
,
a
-\infty -
,
)
,
a

\infty 

上单增, 
在
,
a
a


-

上单减;又(
)
f
a
b
a


-
=
+
>




, (
)
f
a
b
a


=
-
<




,并且 
()
lim
x
f x
\to -\infty 
=-\infty ,
()
lim
x
f x
\to +\infty 
=+\infty ,故方程有3个实根. 
\vspace{6pt}
{\large \bfseries \textcolor{lyred}{第3.5 节 函数的极值与最大值最小值}}

\vspace{4pt}
{\bfseries \textcolor{lyred}{一.函数的极值及其求法}}

\textcolor{lyred}{定义}.设
()
f x 在
(
)
U x
内有定义,若在某
(
)
U x
\circ 
内,
()
(
)
f x
f x
<
,则称
(
)
f x
为 
()
f x 的一个\textcolor{lyred}{极大值};\textcolor{lyred}{极小值}类似.极大值极小值统称为\textcolor{lyred}{极值}(看图). 
\textcolor{lyred}{例}.设
()
f x 连续,
()
lim
n
x
f x
x
\to 
=,讨论
()
f
是否为极值. 
解.
()
f
=
,在0 的两侧邻域内,
()
f x 与
nx 同号,故当n 为奇数时,
()
f
=
不是 
极值,当n 为偶数时,
()
f
=
是极小值. 
\textcolor{lyred}{定理(必要条件)}.设
()
f x 在
0x 处可导,且有极值,则
(
)
f
x
'
=
. 
\textcolor{lyred}{注}.因此,极值点必是驻点或不可导点;反之不对,例如
()
f x
x
=
. 
\textcolor{lyred}{定理(第一充分条件)}.设
()
f x 在
0x 处连续,在
(
)
U x
\circ 
内可导,则 
(1)当在
0x 的左侧
()
f
x
'
>
,右侧
()
f
x
'
<
时,
(
)
f x
为极大值; 
(2)当在
0x 的左侧
()
f
x
'
<
,右侧
()
f
x
'
>
时,
(
)
f x
为极小值; 
(3)当在
0x 的两侧
()
f
x
'
的符号相同时,
(
)
f x
不是极值. 
\textcolor{lyred}{例}.设
()
f x 二阶可导,
()
f '
=
,
()
lim
x
f
x
x
\to 
''
=,证明:
()
f
为极小值. 
证.在0 两侧邻域内,
()
()
f
x
f
x
''
'
>
\Rightarrow 
在
()
U
内单增,故在0 的左侧
()
f
x
'
<
, 
右侧
()
f
x
'
>
,证毕. 
\textcolor{lyred}{例}.求
()
(
)
x
x
f x
x
x
-
=
-
在(
)
0,1 内的极值. 
解.
()
(
)
ln1
x
x
x
f
x
x
x
x
-
'
=
-
-
,驻点
x =
,当
x
<
<
时,
()
f
x
'
<
,当
x >
时, 
()
f
x
'
>
,故
f 
=




为极小值,无极大值. 
\textcolor{lyred}{例}.求
()
(
)(
)
f x
x
x
=
-
+
的极值. 
解.
()
(
)(
)
f
x
x
x
-
'
=
-
+
,驻点
x =,不可导点
x =-,列表如下: 
x        (
)
, 1
-\infty -
-        (
)
1,1
-
        1          (
)
1,+\infty  
()
f
x
'
       +         不可导       -          0            + 
()
f x      单增    极大
(
)1
f -
=
    单减   极小
()
3 4
f
=-
   单增 
\textcolor{lyred}{定理(第二充分条件)}.设
()
f x 在
0x 处二阶可导,且
(
)
f
x
'
=
,则 
(1)当
(
)
f
x
''
<
时,
(
)
f x
为极大值;(2)当
(
)
f
x
''
>
时,
(
)
f x
为极小值. 
\textcolor{lyred}{例}.求
()(
)
f x
x
=
-
+的极值. 
解.
()
(
)
f
x
x x
'
=
-
,驻点
x =
, 1
\pm ,
()
(
)(
)
1 5
f
x
x
x
''
=
-
-
,而 
()
f ''
=
>
,故
()
f
=
是极小值; 
(
)1
f ''\pm 
=
,而在1
\pm 两侧,
()
f
x
'
同号,故
(
)1
f \pm 
不是极值点. 
\textcolor{lyred}{例}.设
()
f
x 在
x
x
=
处二阶可导,若
(
)
(
)
(
)
lim
h
f x
h
f x
f
x
h
\to 
'
+
-
-
=
,证明: 
()
f
x 在
0x 处取到极小值. 
证.
(
)
f
x
'
=
,故
(
)
(
)
(
)
(
)
lim
lim
h
h
f
x
h
f
x
h
f
x
f
x
h
h
\to 
\to 
'
'
'
''
+
+
-
=
=
=
>
,证毕. 
\vspace{4pt}
{\bfseries \textcolor{lyred}{二.最大值最小值问题}}

\textcolor{lyred}{例}.求
()
f x
x
x
=
-
+
在[
]
3,4
-
上的最大值和最小值. 
解.
()
f x 是闭区域上的连续函数,故一定存在最大值和最小值; 
设()(
)
g x
x
x
=
-
+
,则
()
(
)(
)
g
x
x
x
x
'
=
-
+
-
,驻点
x =,2 , 3
2 , 
()1
g
=
, ()
g
=
,
g 
=




, (
)
g -
=
, ()
g
=
,最大值20 ,最小值0 . 
\textcolor{lyred}{例}.铁路AB 距离100km ,工厂C 距A 处20km, AC 垂直于AB ,现在AB 上选D 向 
工厂修筑公路,已知铁路与公路每公里运费之比为3:5,求AD 距离,使得货物从 
工厂运到B 的运费最省. 
解一.设AD
x
=
,
DB
x
=
-
,
CD
x
x
=
+
=
+
,铁路运费每公里3k , 
公路5k ,总运费
(
)
y
k
x
k
x
=
+
+
-
,0
x
\le 
\le 
; 
x
y
k
x


'=
-


+


,唯一驻点
x =
,由(
)
y
k
=
, ()
y
k
=
, 
(
)
y
k
=
,得当
AD =
时,运费最省. 
解二.
x =
是唯一驻点,而
(
)
3 2
k
y
x
''=
>
+
,故为最小值点. 
\textcolor{lyred}{例}.做一个容积为V 的无盖圆柱形桶,底面是铝板,侧面是木板,已知铝板价格是 
木板的5倍,求底面半径使得费用最省. 
解.设底面半径为r ,桶高为
V
h
r
\pi 
=
,则总费用
V
y
k
r
k
r
r
\pi 
\pi 
\pi 
=
\cdots 
+
\cdots 
\cdots 
= 
(
)
V
k
r
r
r
\pi 


+
>




,
dy
V
V
k
r
r
dr
r
\pi 
\pi 


=
-
=
\Rightarrow 
=




唯一驻点,而 
d y
V
k
dr
r
\pi 


=
+
>




,故
V
r
\pi 
=
是极小值点,也是最小值点. 
\textcolor{lyred}{补充练习 }
1.设
()
(
)
lim
x
f
x
x
\to 
'
=-
-
,证明:
()1
f
为极大值. 
证.在
x =两侧邻域内
()
f
x
'
与(
)
x -
同号,即得,证毕. 
2.设
(
)
(
)(
)
n
f
x
f
x
-
'
=
=
=
\Lambda 
,
()(
)
n
f
x
\neq 
,则(1)当n 为偶数时,
(
)
f x
是极值, 
且当
()(
)
n
f
x
<
时为极大值,当
()(
)
n
f
x
>
时为极小值; 
(2)n 为奇数时,
(
)
f x
不是极值. 
证.
()
(
)
()(
)(
)
(
)
()
(
)
(
)
()(
)
lim
!
!
n
n
n
n
n
x
x
f
x
f x
f x
f
x
f x
f x
x
x
o
x
x
n
n
x
x
\to 
-


=
+
-
+
-
\Rightarrow 
=


-
, 
当n 为偶数(奇数)时,在
0x 两侧邻域内
()
(
)
f x
f x
-
同号(不同),证毕. 
3.设a 和b 为两个满足
a
b
+
=的正数,求
(
)
,
m
n
a b
m n >
的最大值. 
解.令x
a
=
,
b
x
=-
,
()
(
)(
)
n
m
f x
x
x
x
=
-
<
<
,则 
()
(
)
(
)
(
)
(
)
n
n
m
m
m
f
x
x
x
m
x
nx
m
n x
x
x
m
n
-
-
-
-


'
=
-
-
-
=
+
-
-






+


,唯一驻点 
m
x
m
n
=
+
,当
m
x
m
n
>
+
时,
()
f
x
'
<
,当
m
x
m
n
<
+
时,
()
f
x
'
>
,故最大值为 
m
n
m
m
n
f
m
n
m
n
m
n





=





+
+
+





. 
4.设a 和b 为两个满足
ab =的正数,求
(
)
,
m
n
a
b
m n
+
>
的最小值. 
解.令x
a
=
,
b
x
=
,
()
(
)
m
n
f x
x
x
x
-
=
+
>
,则
()
m
n
f
x
mx
nx
-
--
'
=
-
= 
n
m n
n
mx
x
m
--
+


-




,唯一驻点
m n
n
x
m
+


=



,而当
x
x
<
时,
()
f
x
'
<
,当
x
x
>
时, 
()
f
x
'
>
,故最小值为
(
)
m
n
m n
m n
n
m
f x
m
n
+
+




=
+








. 
\vspace{6pt}
{\large \bfseries \textcolor{lyred}{第3.6 节 函数图形的描绘}}

借助一阶导数的符号,可以确定函数图形的上升下降区间,借助二阶导数的符号, 
可以确定函数图形的凹凸区间. 
\textcolor{lyred}{例}.描绘
y
x
x
x
=
-
-
+的图形. 
解.(1)定义域为(
)
,
-\infty +\infty ; 
(2)
(
)(
)
y
x
x
x
'=
+
-
=
\Rightarrow 
=-
,1;
y
x
x


''=
-
=
\Rightarrow 
=




; 
(3)列表如下: 
x    
,


-\infty -




-
1 1
,
3 3


-




1     1 ,1






      1    (
)
1,+\infty  
y'       +        0        -       -       -        0     + 
y''       -        -        -       0       +       +     + 
图      升凸      极大      降凸    拐点     降凹     极小   升凹 
(4) lim
x
y
\to -\infty 
=-\infty , lim
x
y
\to +\infty 
=+\infty ; 
(5)
1 32
,
3 27


-




, 1 16
,
3 27






,(
)
1,0 ,补充(
)
1,0
-
,(
)
0,1 . 
\textcolor{lyred}{例}.描绘
x
y
e
\pi 
-
=
的图形. 
解.(1)定义域为(
)
,
-\infty +\infty ,偶函数,且
y >
; 
(2)
x
y
xe
x
\pi 
-
'=-
=
\Rightarrow 
=
;
(
)
x
y
e
x
x
\pi 
-
''=
-
=
\Rightarrow 
=\pm ; 
(3)列表如下: 
x     (
)
, 1
-\infty -
-    (
)
1,0
-
     0     (
)
0,1     1    (
)
1,+\infty  
y'       +       +       +       0      -      -      - 
y''       +       0       -       -      -      0      + 
图      升凹     拐点     升凸    极大    降凸    拐点    降凹 
(4)lim
x
y
\to \infty 
=
,水平渐近线
y =
; 
(5)
0,
2\pi 






,
1,
2 e
\pi 


\pm 




. 
\textcolor{lyred}{例}.描绘
x
y
xe
=
的图形. 
解.(1)定义域(
)
(
)
,0
0,
-\infty 
\cup 
+\infty ,间断点
x =
; 
(2)
x
x
y
e
x
x
-
'=
=
\Rightarrow 
=;
x
y
e
x
''=
; 
(3)列表如下: 
x        (
)
,0
-\infty 
      0      (
)
0,1      1      (
)
1,+\infty  
y'          +        \times        -        0        + 
y''          -        \times        +        +       + 
图         升凸       \times       降凹      极小     升凹 
(4)
lim
x
y
-
\to 
=
,
lim
x
y
+
\to 
=+\infty ,铅直渐近线
x =
, lim
x
y
\to -\infty 
=-\infty , lim
x
y
\to +\infty 
=+\infty ; 
(5)(
)
1,e . 
\textcolor{lyred}{例}.描绘
(
)
x
y
x
=+
+
的图形. 
解.(1)定义域(
)
(
)
, 3
3,
-\infty -
\cup -
+\infty ,间断点为
x =-; 
(2)
(
)
(
)
36 3
x
y
x
x
-
'=
=
\Rightarrow 
=
+
;
(
)
(
)
x
y
x
x
-
''=
=
\Rightarrow 
=
+
; 
(3)列表如下: 
x     (
)
, 3
-\infty -
-    (
)
3,3
-
    3     (
)
3,6     6     (
)
6,+\infty  
y'        -       \times       +      0       -      -       - 
y''        -       \times       -      -       -      0       + 
图       降凸      \times      升凸    极大    降凸    拐点     降凹 
(4)
lim
x
y
\to -
=-\infty ,铅直渐近线
x =-,lim
x
y
\to \infty 
=,水平渐近线
y =; 
(5)(
)
3,4 ,
6, 3






,补充(
)
0,1 . 
\textcolor{lyred}{补充练习 }
1.描绘
(
)
x
y
x
=
-
的图形. 
解.(1)定义域(
)
(
)
,1
1,
-\infty 
\cup 
+\infty ,间断点为
x =; 
(2)
(
)
(
)
x
x
y
x
x
-
'=
=
\Rightarrow 
=
-
,3;
(
)
x
y
x
x
''=
=
\Rightarrow 
=
-
,1; 
(3)列表如下: 
x     (
)
,0
-\infty 
     0     (
)
0,1    1     (
)
1,3     3     (
)
3,+\infty  
y'       +       0      +      \times       -      0       + 
y''       -       0      +      \times       +      +       + 
图      升凸     拐点    升凹     \times      降凹    极小    升凹 
(4)
lim
x
y
\to 
=+\infty ,铅直渐近线
x =, lim
x
y
\to -\infty 
=-\infty , lim
x
y
\to +\infty 
=+\infty ; 
(5)(
)
0,0 ,
3, 4






. 
2.描绘
(
)
x
y
x
-
=
-
的图形. 
解.(1)定义域(
)
(
)
,1
1,
-\infty 
\cup 
+\infty ,间断点为
x =; 
(2)
(
)
x
y
x
x
-
'=
=
\Rightarrow 
=
-
;
(
)
(
)
2 2
x
y
x
x
+
''=
=
\Rightarrow 
=-
-
; 
(3)列表如下: 
x    
,


-\infty -




-
1 ,0


-




    0    (
)
0,1   1    (
)
1,+\infty  
y'       -       -       -       0     +    \times       - 
y''      -        0       +      +     +    \times       + 
图      降凸     拐点    降凹     极小   升凹   \times      降凹 
(4)
lim
x
y
\to 
=+\infty ,铅直渐近线
x =,lim
x
y
\to \infty 
=
,水平渐近线
y =
; 
(5)
,


-
-




,(
)
0, 1
-
. 
\vspace{6pt}
{\large \bfseries \textcolor{lyred}{第3.7 节 曲率}}

\vspace{4pt}
{\bfseries \textcolor{lyred}{一.弧微分}}

\textcolor{lyred}{定义}.设L :
()
y
f x
=
为\textcolor{lyred}{光滑曲线},取点
(
)
,
M
x
y
L
\in 
,以x 增大的方向为曲线的 
正方向,
(
)
,
M x y
L
\forall 
\in 
,令 
()
\dots 
\dots 
,
,
M M
x
x
s x
M M
x
x
-
<

=
\ge 

,称为\textcolor{lyred}{有向弧段的值},或\textcolor{lyred}{长度},简称为\textcolor{lyred}{弧},也记为\dots 
M M . 
\textcolor{lyred}{弧微分公式}.
()
ds
f
x
dx
'
=
+
.\textcolor{lyred}{ }
\vspace{4pt}
{\bfseries \textcolor{lyred}{二.曲率及其计算公式}}

\textcolor{lyred}{定义}.设L 为光滑曲线,基点
M ,
(
)
,
M x y
L
\in 
对应弧长()
s x ,切线倾角为
()
x
\alpha 
, 
(
)
,
M
x
x y
y
L
\Delta 
\Delta 
'
+
+
\in 
,对应弧长(
)
s x
x
\Delta 
+
,切线倾角为(
)
x
x
\alpha 
+\Delta 
,则 
(
)
()
(
)
()
x
x
x
K
s
s x
x
s x
\alpha 
\Delta 
\alpha 
\Delta \alpha 
\Delta 
\Delta 
+
-
=
=
+
-
,称为弧段\dots 
MM '的\textcolor{lyred}{平均曲率},而M 处的\textcolor{lyred}{曲率}为\textcolor{lyred}{ }
lim
M
M
M
d
K
K
ds
\alpha 
'\to 
=
=
,即切线方向角对弧长的变化率. 
\textcolor{lyred}{例}.对于直线,任意两点之间的
\Delta \alpha =
,故
K
s
\Delta \alpha 
\Delta 
=
=
,而
K =
. 
\textcolor{lyred}{例}.对于半径为r 的圆,
s
r
r
\Delta \alpha 
\Delta 
\pi 
\Delta \alpha 
\pi 
=
\cdots 
=
,故
K
s
r
\Delta \alpha 
\Delta 
=
=
,而
K
r
=
. 
\textcolor{lyred}{公式}.光滑曲线
()
:
L y
f x
=
在(
)
,x y 处的曲率为
(
)
y
K
y
''
=
'
+
.\textcolor{lyred}{ }
\textcolor{lyred}{推论}.设光滑曲线
()
()
:
x
t
L
y
t
\varphi 
\zeta 
=

=

,若
\varphi 
\zeta 
'
'
+
\neq 
,则
K
\varphi \zeta 
\varphi \zeta 
\varphi 
\zeta 
'''
'''
-
=
'
'


+


. 
\textcolor{lyred}{例}.求双曲线
xy =在(
)
1,1 处的曲率. 
解.
y
x
'=-
,
y
x
''=
,故
(
)
K =
=


+-


. 
\textcolor{lyred}{例}.求双曲线
y
x
x
=
-
在(
)
2,0 处的曲率. 
解.
dy
dy
x
y
x
dx
dx
y
-
\cdots 
=
-
\Rightarrow 
=
,故在(
)
2,0 处, dy
dx 不存在,而
dx
dy =
,
d x
dy
=
, 
故
(
)
(
)
1 0
x
K
x
''
=
=
=
+
'
+
. 
\textcolor{lyred}{例}.求抛物线
y
ax
bx
c
=
+
+上曲率的最大值. 
解.
(
)
a
K
ax
b
=


+
+


,当2
ax
b
+
=
,即
a
b
x
-
=
时,
max
K
a
=
. 
\textcolor{lyred}{例}.求摆线
(
)
(
)(
)
sin
1 cos
x
a t
t
t
y
a
t
\pi 
=
-

<<

=
-

的最小曲率. 
解.
(
)
(
)
8 1 cos
x y
x y
K
a
t
x
y
'''
'''
-
=
=
-
'
'
+
,当t
\pi 
=
时,
min
K
a
=
. 
\vspace{4pt}
{\bfseries \textcolor{lyred}{三.曲率半径与曲率圆}}

设曲线L 在M 处的曲率
K \neq 
,则在它凹侧的法线上取一点D ,使得
DM
K
=
, 
以D 为圆心, 1
K 为半径作圆,称为L 在M 处的\textcolor{lyred}{曲率圆}, D 为\textcolor{lyred}{曲率中心},
K
\rho =
为 
\textcolor{lyred}{曲率半径}. 
\textcolor{lyred}{例}.设5吨的汽车以每秒6 米的速度在跨度10米,拱高0.25米的抛物线型拱桥上 
行驶,求它越过桥顶时对桥面的压力. 
解.设桥面的方程为
y
ax
=
,代入
x =
,
0.25
y =-
,得
0.01
a =-
,桥顶处曲率 
0.02
K
a
=
=
,曲率半径
R
K
=
=
米,于是向心力为
mv
mg
F
R
-
=
,得 
(
)
45400 N
mv
F
mg
R
=
-
=
. 